
\begin{proposition}
\label{XII.5.3}
Sean $\othercal{X}$ un espacio analítico normal y $\othercal{Y}$ un
subconjunto analítico cerrado tal que
$\othercal{U}=\othercal{X}-\othercal{Y}$ sea denso en~$\othercal{X}$.
Entonces el funtor que a cada recubrimiento finito
normal~\eqref{XII.5.0} $\othercal{X}'$ de~$\othercal{X}$ le asigna su
restricción a~$\othercal{U}$ es plenamente fiel.
\end{proposition}

Sean $\othercal{X}'$
\ifthenelse{\boolean{orig}}
{y}
{y}%
$\othercal{X}''$ dos recubrimientos finitos normales de~$\othercal{X}$.
Debemos demostrar que la aplicación canónica
$$
\Hom_{\othercal{X}}(\othercal{X}',\othercal{X}'') \to
\Hom_{\othercal{U}}(\othercal{X}'|\othercal{U},\othercal{X}''|\othercal{U})
$$
es biyectiva. Sean $u$, $v$ dos $\othercal{X}$-morfismos de
$\othercal{X}'$ en~$\othercal{X}''$ cuyas restricciones a~$\othercal{U}$
son iguales y demostremos que $u=v$. Los morfismos $u$ y
\marginpar{338}
$v$ coinciden sobre el abierto denso
$\othercal{U}\times_{\othercal{X}}\othercal{X}'$ y, por tanto, sobre los
espacios topológicos subyacentes. Según
\hbox{\cite[\no 19~\S 4~cor\ptbl 5]{XII.4}}, esto demuestra que $u=v$.

Sea ahora $u$ un $\othercal{U}$-morfismo de
$\othercal{X}'|\othercal{U}$ en~$\othercal{X}''|\othercal{U}$ y
demostremos que se prolonga a todo~$\othercal{X}'$. Podemos suponer que
$\othercal{X}'$ es regular. En efecto, como $\othercal{X}'$ es normal,
se puede encontrar un abierto $\othercal{V}$ de~$\othercal{X}$ cuyo
complementario sea un subconjunto analítico de codimensión $\geq2$, tal
que $\othercal{X}'\times_{\othercal{X}}\othercal{V}=\othercal{V}'$ sea
regular. Sea
$\othercal{V}''=\othercal{X}''\times_{\othercal{X}}\othercal{V}$ y
supongamos demostrada la proposición para~$\othercal{V}$. Consideremos el
diagrama conmutativo
\begingroup
\shorthandoff{>}% Spanish babel otherwise treats the Xy-pic arrow tip as a shorthand.
$$
\xymatrix{ \othercal{V}' \ar[dd]_{g'}\ar[rd]\ar[rrr]^{i'} && &\othercal{X}'
\ar[dd]^{\hbox{\raise10mm\hbox{$f'$}}} \cr &\othercal{V}'' \ar[ld]_{g''\mkern-10mu}\ar[rrr]^-{i''~~} & &
&\othercal{X}'' \ar[ld]_{f''\mkern-12mu} \cr \othercal{V} \ar[rrr]^{i} && &\othercal{X} \cr
}
$$
\endgroup
A~$u$ se le asocia un morfismo de
$\cal{O}_{\othercal{V}}$-álgebras
$g''_*\cal{O}_{\othercal{V}''}\to g'_*\cal{O}_{\othercal{V}'}$, de donde
se deduce un morfismo
$$
i_*g''_*\cal{O}_{\othercal{V}''}\to
i_*g'_*\cal{O}_{\othercal{V}'}\quoi.
$$
Teniendo en cuenta los isomorfismos
$i'_*\cal{O}_{\othercal{V}'}\simeq\cal{O}_{\othercal{X}'}$ e
$i''_*\cal{O}_{\othercal{V}''}\simeq\cal{O}_{\othercal{X}''}$
\cite[\no 3~prop\ptbl 4]{XII.11}, se deduce un morfismo de
$\cal{O}_{\othercal{X}}$-álgebras
$$
f''_*\cal{O}_{\othercal{X}''}\to
f'_*\cal{O}_{\othercal{X}'}\quoi,
$$
de donde se obtiene el morfismo buscado
$\othercal{X}'\to\othercal{X}''$.

Supongamos
\marginpar{339}
en adelante que $\othercal{X}'$ es regular. Sean
$\othercal{U}'=\othercal{U}\times_{\othercal{X}}\othercal{X}'$ y
$\othercal{Y}'=\othercal{X}'-\othercal{U}'$. Consideremos
$\othercal{Y}'$ como subespacio analítico reducido de~$\othercal{X}'$;
si $\othercal{Y}'_1$ es el lugar singular de~$\othercal{Y}'$, se tiene
$\dim\othercal{Y}'_1<\dim\othercal{Y}'$
\cite[\no 20~D\ptbl Th\ptbl 3]{XII.4}. Por recurrencia sobre la dimensión
de~$\othercal{Y}'$, podemos suponer, pues, que $\othercal{Y}'$ es liso.
Como basta prolongar $u$ a un entorno abierto de cada punto de
$\othercal{Y}'$, por el teorema de las funciones implícitas podemos
suponer que $\othercal{X}'$ es una bola de un espacio afín
$\othercal{E}^n$ y que $\othercal{Y}'$ es el cerrado definido por la
anulación de las primeras $p$ funciones coordenadas
$z_1,\dots,z_p$, con $0\leq p\leq n$.

A~$u$ le asociamos una sección $s$ de
$p\colon\othercal{X}'\times_{\othercal{X}}\othercal{X}''\to\othercal{X}'$
sobre~$\othercal{U}'$; restringiendo $\othercal{X}'$ si es necesario,
podemos suponer que
$p_*(\cal{O}_{\othercal{X}'\times_{\othercal{X}}\othercal{X}''})$ está
engendrado por elementos $x_1,\dots,x_q$ de
$\Gamma(\othercal{X}',p_*\cal{O}_{\othercal{X}'\times_{\othercal{X}}\othercal{X}''})$;
sean $u_1,\dots,u_q\in\Gamma(\othercal{U}',\cal{O}_{\othercal{X}'})$ las
imágenes por~$s$ de~${x_1}|\othercal{U}',\dots,{x_q}|\othercal{U}'$.
Decir que $s$ se prolonga a~$\othercal{X}'$ equivale a decir que
$u_1,\dots,u_q$ se prolongan en secciones de
$\Gamma(\othercal{X}',\cal{O}_{\othercal{X}'})$. Pero como
\ifthenelse{\boolean{orig}}
{$f$}
{$p$}%
~es finito, cada $u_i$ es una serie de Laurent en
$z_1,\dots,z_p$, con coeficientes que son series de potencias en
$z_{p{+}1},\dots,z_n$, y satisface relaciones de dependencia integral.
Se deduce que $u_i$ es acotado, luego es una serie de potencias en
$z_1,\dots,z_n$ y, por consiguiente, se prolonga a~$\othercal{X}'$.

Cabe preguntarse si el funtor introducido en~\Ref{XII.5.3} es una
equivalencia de categorías. El teorema de
\textsc{Grauert}--\textsc{Remmert} \cite{XII.6} da una respuesta a esta
pregunta; a continuación presentamos una demostración que utiliza la
resolución de singularidades. También se habría podido emplear el teorema
de \textsc{Grauert}--\textsc{Remmert} para demostrar~\Ref{XII.5.1}; así
se hacía antes de disponer de~\cite{XII.8}.

\ifthenelse{\boolean{orig}}
{}
{\enlargethispage{.5cm}}
\begin{theoreme}[Teorema de \textsc{Grauert}--\textsc{Remmert}]
\label{XII.5.4}
Sean
\marginpar{340}
\index{Grauert-Remmert (teorema de)|hyperpage}%
\index{teorema de Grauert-Remmert|hyperpage}%
$\othercal{X}$ un espacio analítico normal y $\othercal{Y}$ un
subconjunto analítico cerrado tal que
$\othercal{U}=\othercal{X}-\othercal{Y}$ sea denso en~$\othercal{X}$.
Sea $\othercal{U}'$ un recubrimiento finito normal de~$\othercal{U}$;
supongamos que existe un subconjunto analítico cerrado de interior vacío
$\othercal{S}$ de~$\othercal{X}$ tal que la restricción de
$\othercal{U}'$ a~$\othercal{U}-\othercal{U}\cap\othercal{S}$ sea étale.
Entonces existe un recubrimiento finito normal $\othercal{X}'$ de
$\othercal{X}$ que prolonga $\othercal{U}'$, y $\othercal{X}'$ es único
salvo isomorfismo.
\end{theoreme}

La unicidad resulta de~\Ref{XII.5.3}. El problema de prolongar
$\othercal{U}'$ es, por tanto, local sobre~$\othercal{X}$. Podemos
suponer que $\othercal{U}$ es regular y que $\othercal{U}'$ es étale
sobre~$\othercal{U}$. En efecto, el conjunto de los puntos regulares de
$\othercal{U}$ es un abierto $\othercal{V}$ denso en~$\othercal{X}$ cuyo
complementario es un subconjunto analítico
\hbox{\cite[\no20~D th\ptbl 2]{XII.4}}, y basta sustituir
$\othercal{U}$ por el abierto
$\othercal{V}-\othercal{V}\cap\othercal{S}$.

Sea $y$ un punto de~$\othercal{X}-\othercal{U}$ y demostremos que se
puede prolongar $\othercal{U}'$ a un entorno de~$y$. Restringiendo
$\othercal{X}$ a un entorno abierto de~$y$, resulta del teorema de
resolución de singularidades \cite{XII.8} que se pueden encontrar un
espacio analítico regular $\othercal{X}_1$ y un morfismo proyectivo
$f\colon\othercal{X}_1\to\othercal{X}$ que induce por restricción a
$\othercal{U}$ un isomorfismo
$\othercal{U}_1=f^{-1}(\othercal{U})\simeq\othercal{U}$, tales que
$\othercal{U}_1$ sea el complementario en~$\othercal{X}_1$ de un divisor
con cruzamientos normales. Demostremos que $\othercal{U}'$ se prolonga
en un recubrimiento finito normal de~$\othercal{X}_1$. Como la cuestión
es local sobre~$\othercal{X}_1$, podemos suponer que $\othercal{X}_1$ es
una bola de un espacio afín $\othercal{E}^n$ y que
$\othercal{X}_1-\othercal{U}_1$ está definido por la anulación de las
primeras $p$ funciones coordenadas
\ifthenelse{\boolean{orig}}
{$z_1,\dots,z_q$,}
{$z_1,\dots,z_p$,}
con $0\leq p\leq n$. El recubrimiento étale $\othercal{U}'$ de
$\othercal{U}_1$ es cociente de un recubrimiento de la forma
$$
\othercal{U}_2=\othercal{U}_1[T_1,\dots,T_p]/\left(T_1^{n_1}-z_1,\dots,
T_p^{n_p}-z_p\right)
$$
por un subgrupo $H$ del grupo de Galois de~$\othercal{U}_2$. El
recubrimiento $\othercal{U}_2$ se prolonga
\marginpar{341}
en el recubrimiento
$$
\othercal{X}_2=\othercal{X}_1[T_1,\dots,T_p]/\left(T_1^{n_1}-z_1,\dots,
T_p^{n_p}-z_p\right)
$$
de~$\othercal{X}_1$, sobre el cual actúa $H$, y
$\othercal{X}_2/H$ prolonga $\othercal{U}'$ a~$\othercal{X}_1$.

Denotemos por~$\othercal{X}_1'$ el recubrimiento finito normal de
$\othercal{X}_1$ que prolonga $\othercal{U}'$ y por~$\othercal{F}_1$ la
$\cal{O}_{\othercal{X}_1}$-álgebra coherente
\ifthenelse{\boolean{orig}}
{definida por~$\othercal{F}_1$}
{que define~$\othercal{X}_1'$}%
. Según el teorema de finitud de
\textsc{Grauert}--\textsc{Remmert}
\hbox{\cite[\no15 th\ptbl 1.1]{XII.4}}, $f_*\othercal{F}_1$ es una
$\cal{O}_{\othercal{X}}$-álgebra coherente. Le corresponde, pues, un
recubrimiento finito $\othercal{X}'$ de~$\othercal{X}$, que además es
normal porque $\othercal{X}_1'$ lo es, y $\othercal{X}'$ es la
prolongación buscada de~$\othercal{U}'$.

\begin{remarque}
\label{XII.5.5}
En el enunciado~\Ref{XII.5.4} no se puede suprimir la hipótesis sobre el
lugar de los puntos en los que el morfismo
$\othercal{U}'\to\othercal{U}$ no es étale. Sea, por ejemplo,
$\othercal{X}$ el disco unidad del plano complejo, $\othercal{U}$ el
complementario del origen en~$\othercal{X}$ y
$\othercal{U}'=\othercal{U}[T]/(T^2-\sin 1/z)$, donde $z$ es la función
coordenada sobre~$\othercal{X}$. Entonces $\othercal{U}'$ es un
recubrimiento finito normal de~$\othercal{U}$ que no se prolonga a
$\othercal{X}$. Supongamos, en efecto, que $\othercal{U}'$ se prolonga
en un recubrimiento finito $\othercal{X}'$ de~$\othercal{X}$; el lugar de
los puntos de~$\othercal{X}$ en los que el morfismo
$\othercal{X}'\to\othercal{X}$ no es étale sería entonces un cerrado
analítico que contendría todos los puntos $z$ tales que
$\sin 1/z=0$, lo cual es absurdo.

Sin embargo, se puede suprimir la hipótesis sobre el lugar singular del
morfismo $\othercal{U}'\to\othercal{U}$ cuando
$\codim(\othercal{X}-\othercal{U},\othercal{X})\geq2$. Podemos suponer,
en efecto, que $\othercal{U}$ es regular. El lugar de los puntos de
$\othercal{U}$ en los que $\othercal{U}'\to\othercal{U}$ no es étale es
un divisor de~$\othercal{U}$, y del teorema de
\textsc{Remmert}--\textsc{Stein} \cite[th\ptbl 3]{XII.9} resulta que es
\marginpar{342}
la traza sobre~$\othercal{U}$ de un divisor de~$\othercal{X}$. Ahora
bien, en este caso, si $\othercal{A}$ es una
$\cal{O}_{\othercal{U}}$-álgebra coherente tal que
$\othercal{U}'=\Spec\an(\othercal{A})$ y si
$i:\othercal{U}\to\othercal{X}$ es el morfismo canónico, basta
tomar $\othercal{X}'=\Spec\an(i_*\othercal{A})$; se sabe, en efecto, que
$i_*\othercal{A}$ es
\ifthenelse{\boolean{orig}}
{coherentes}
{coherente}
\cite[\no 1 th\ptbl 1]{XII.11}.
\end{remarque}

\ifthenelse{\boolean{orig}}
{}
{\begin{remarqueMR}
\label{XII.5.6}
Existen grupos no triviales finitamente presentados $G$ que no poseen
subgrupos de índice finito distintos de~$G$, por ejemplo el grupo de
G\ptbl Higman (\cf J-P\ptbl Serre, \emph{Arbres et amalgames}, Astérisque
N$^\circ$46, prop\ptbl 6, cap\ptbl I, \S 1). Por consiguiente, si un grupo
de este tipo se realiza como grupo fundamental topológico de un esquema
$V$ sobre~$\CC$, digamos liso y proyectivo, el espacio topológico
$V_{\mathrm{top}}$ subyacente a~$V$ no es simplemente conexo, pero $V$
es algebraicamente simplemente conexo. En la actualidad no se conoce
ningún $V$ de este tipo. Señalemos, no obstante, que D\ptbl Toledo ha
construido esquemas lisos y proyectivos $V$ sobre~$\CC$ cuyo grupo
fundamental topológico no es Hausdorff para la topología de los
subgrupos de índice finito (el morfismo natural
$\pi_1(V_{\mathrm{top}})\to\pi_1(V_{\alg})$ no es inyectivo).
[D\ptbl Toledo, Projective varieties with non-residually finite
fundamental group, Publ.\ Math.\ IHES \textbf{77} (1993),
p\ptbl 103--119].
\end{remarqueMR}}
